\input fontmac
\input mathmac

\input epsf
\widemargins
\bookheader{\tenpoint\sc marcel goh}{\tenpoint\sc math {\oldstyle 140} fall {\oldstyle 2026}}

\def\datestamp#1#2 {\llap{{\oldstyle #1}{\rm.}{\sc #2}\hskip12pt}}
\def\bar{\overline}
\def\symdiff{\mathrel{\triangle}}
\def\To{\Rightarrow}
\def\Gets{\Leftarrow}
\font\titlefont=cmbx12 at 14 pt
\def\cont#1#2{#1\hskip-2pt\leaders\hbox to 8pt{\hss.\hss}\hfil\hbox to 2em{\hss#2}}
\def\divides{\mid}
\def\percent{\mathrel{\%}}
\def\norm#1{|\!|#1|\!|}
\def\dist{\op{\rm dist}}

\maketitle{MATH 140 Fall 2026}{notes by}{Marcel Goh}{}

\floattext 4.5 \ninebf A note on these notes.
\ninepoint This document will periodically be updated with the material that we cover.
The datestamps in the left margin indicate when the notes from each day start.
Subsections labelled with a $*$ are optional.
These notes should {\it not} be used as a substitute for the textbook or for attending class.
(For instance, on the blackboard I will often draw graphs or diagrams, and these for the most part
will be omitted from the notes.)
The best way to do well
in this class is to do as many exercises from Stewart's textbook as you can stomach.
This document may contain errors, typographical or mathematical.
Please email me if you find any.

\vskip36pt

\vskip \medskipamount \hbox to\hsize
{\relax \unhcopy \strutbox \cont
{{\tenssbx I. Foundations \kern 3pt}}{2}}
\smallskip
\hbox to\hsize {\relax \unhcopy \strutbox \cont {\hskip10pt 1.\kern .5em Rates of change}{3}}
\hbox to\hsize {\relax \unhcopy \strutbox \cont {\hskip10pt 2.\kern .5em Limits}{5}}
\medskip

\vfill\eject
\begingroup\headline{\hfil}\footline{\hfil}
\vglue160pt
\centerline{\titlefont I. FOUNDATIONS}
\vskip220pt
\bigskip
\begingroup\obeylines\eightssi
\hfill Que si cependant la Geometrie a to\^ujours
\hfill quelque obscurit\'e essentiel, qu'on ne puisse dissiper,
\hfill \& ce fera uniquement, \`a ce que je crois, du c\^ot\'e de l'Infini,
\hfill c'est que de ce c\^ot\'e-l\`a la Geometrie tient \`a la Phisique,
\hfill \`a la nature intime des Corps que nous connoissons peu,
\hfill \& peut-\^etre aussi \`a une Metaphisique trop \'elev\'ee,
\hfill dont il ne nous est permis que d'appercevoir quelques rayons.
\eightss
\smallskip
\hfill --- BERNARD LE BOVIER DE FONTENELLE, {\eightssi \'Elements de la Geometrie de l'Infini} (1727)
\endgroup%\obeylines
\bigskip\goodbreak
\vfill\eject
\endgroup%blank headline and footline


\advsect Rates of change

\datestamp{31}{viii}
This set of notes assumes that the reader is familiar with all the pre-calculus concepts that one learns
in the later years of secondary school:~functions, trigonometry, exponential functions, logarithms, etc.
However, given that the very definition of what differential calculus is requires us to talk about functions
and their inputs, it is perhaps prudent to begin with a refresher on real-valued functions with real-valued
outputs.

A {\it function} is a collection of ordered pairs $(x,y)$ with the rule that if $(x,y)$
and $(x,z)$ are both in the collection, then $y=z$. We can view $x$ as the {\it input} to the function,
and $y$ as the {\it output} of the function; under this view, the rule above is simply that the function
never outputs two different values when given the same input. People like to use the letter
$f$ (or $g$ or $h$, but really we can use any letter) to denote a function,
and $f(x)$ to denote its value evaluated at a given input $x$. However, many particular functions are
difficult to write without involving the input variable $x$ in the expression,
and in these cases we'll write $f(x)$ to mean the function as well.

Here are some examples of functions:
\medskip
\item{a)} The function $f(x) = 3x$, defined on all real numbers, that outputs triple its input value.
\smallskip
\item{b)} The function that assigns to each real number $x$ its square $x^2$.
\smallskip
\item{c)} The function $y=f(x)$ that equals $0$ whenever $x$ is a negative real number, and equals $x$
whenever $x\ge 0$.
\smallskip
\item{d)} The function that takes as input any nonzero real number $x$ and outputs its reciprocal $1/x$.
\smallskip
\item{e)} The function $f(x)$, defined on all real numbers, given by
$$f(x) = \cases{ 1, & if $x$ is rational;\cr 0, & if $x$ is irrational.}$$
\medskip
Note that in each case above, we were careful to specify what the possible inputs to the function were. Going back
to the definition of a function as a collection of ordered pairs $(x,y)$, the set of all $x$ such
that $(x,y)$ is in the function is called the {\it domain} of $f$; this is precisely the set of all valid
inputs to the function $f$. (The set of all $y$ such that $(x,y)$
is in $f$ is called the {\it range} of $f$, but this definition won't be as important to us.)

The title of this section is ``rates of change,'' but determining the rate at which a given function changes
is actually what this {\it entire class}
will essentially be about. We want to study how much the output of a function changes when you change its input.

For example, consider a car moving from Montreal to Quebec City. As time goes on, it gets further away from
Montreal and closer to Quebec City. So we can model the trip as a function $d(t)$, where for each positive
real number $t$ we let $d(t)$ equal the distance of the car from Montreal (in kilometres),
exactly $t$ minutes since it
set off. The rate of change of this function at any given point in time $t$ is the car's {\it velocity} at that
point.

How do we calculate velocity? Well, we know that it is measured in units such as ``kilometres per hour''
or ``metres per second,'' so it is not too hard to see that velocity is calculated by dividing distance
over time.

Let's go back to the car example. Suppose your car is equipped with a very good GPS that can measure your
distance from Montreal; however, it does not have a working speedometer. Your friend (who is driving) asks
you, ``How fast am I driving?'' You have taken high school physics, so you know how to solve this problem.
You get out a stopwatch and start it exactly when the GPS measures your distance at $135.2$ kilometres
from Montreal. After one
minute on the stopwatch, the GPS reads $137.4$, so you report to your friend that she is driving
$137.4-135.2 = 2.2$ kilometres per minute, or $132$ kilometres per hour.

But your friend is not satisfied. ``You told me my average speed over the course of a minute, not my speed
at the moment I asked you the question. During that minute, I had to speed up to overtake a U-Haul and then
slow down because someone cut me off in the fast lane.''

So you try and do better. You get out the stopwatch again, and when you start it the GPS reads $137.7$.
You wait ten seconds, then observe that the GPS reads $138.1$. So you report your friend's new speed at
$138.1-137.7 = 0.4$ kilometres every ten seconds, or $2.4$ kilometres per minute, or $144$ kilometres per
hour.

Your friend is a little more satisfied, but wonders if you can do better. So you try one last time. You
start the stopwatch when the GPS reads $138.5$, wait exactly one second, and then check the GPS again.
To everyone's surprise and disappointment, it still reads $138.5$, because it only measures distance to the
tenth part of a kilometre. So the calculation would report
that the car has moved $0$ kilometres in that one second, that is, it's going $0$ kilometres per hour.

Obviously in this example you were hamstrung by the precision of your measuring instruments, but it is still
worth pondering the following point. Your friend was asking for her speed at some particular instant.
However, during any particular instant she does not move, nor does time move. So any na{\"\i}ve calculation
as to her speed would go along the lines of $0/0$, which is an undefined value.

Yet it is undeniable that something's ``speed at a particular instant'' is a concept that holds meaning
in everyday life. If you are in a plane and your friend is in a car below you, then it's certainly
not nonsensical to say that you are moving faster than your friend at that moment, even though technically it
could also be said that both of you have moved ``$0$ kilometres over $0$ hours.''

The way we make this precise is suggested by your technique in the first scenario. To avoid dividing by zero,
we let time elapse just a tiny bit, then measure how much the object's position has changed in that minute
interval, and perform the division. Assuming we have arbitrarily sophisticated measuring instruments,
we can keep taking smaller and smaller intervals, and arrive at something closer and closer to the
``instantaneous velocity.'' But it is worth stating that ``instantaneous velocity'' is a theoretical concept,
and not one that can be measured exactly. So when you are driving in a real car, and your speedometer
reads some value, say, $127$ kilometres per hour, what it's really doing is measuring your car's
change in distance over a very small period of time, and reporting to you the quotient (not unlike what
you were doing in that first scenario!).

\medskip\boldlabel Secants and tangents.
The average rate of change of a function $f(x)$
between the input $x$ and the input $x+h$ (for
some small $h>0$) is
the total amount it changes, namely $f(x+h) - f(x)$, divided by the amount the input changed, namely
$(x+h)-x = h$.
If you graph the function $f(x)$ on the $(x,y)$ plane, then geometrically
what you are doing is drawing a line between the points
$\bigl( x, f(x)\bigr)$ and $\bigl( x+h, f(x+h)\bigr)$ and calculating its {\it slope} (the ``rise over
run''). Such a line cutting through the function's graph is called a {\it secant line}, because
it {\it cuts} through the function, at the very least at those two aforementioned points.

By taking $h$ smaller and smaller (but never zero!) we continue to obtain secant lines, but the distance
between the two points $\bigl( x, f(x)\bigr)$ and $\bigl( x+h, f(x+h)\bigr)$ will get smaller and smaller.
You can imagine these points getting so close that it's as if the {\sl line touches the function exactly
once}, at the point $\bigl(x, f(x)\bigr)$ and nowhere else. This line, if it exists, will be called
the {\it tangent line}, but we're not ready to precisely define it just yet.

All we note at this point is that if we admit the concept of the ``instantaneous rate of change,'' and
allow that it can be approximated by taking the the average rate of change over smaller and smaller
intervals, then the slope of this tangent line (if it exists!) is exactly the instantaneous rate of change
that we seek. So it seems that the slope of the tangent line should be the value of the quotient
$$ {f(x+h) - f(x) \over (x+h) - x} = {f(x+h) - f(x) \over h}$$
when $h=0$, if only we could plug in $h=0$. But we cannot, since division by zero is undefined.

What we have learned from this preliminary section is that to understand rates of change,
we shall have to develop a precise concept of what it means to take a function
close to some input at which it is not necessarily defined. This is the notion of a {\it limit}.

\advsect Limits

To understand the definition of a limit, it might be helpful to think about it as a generalisation of the concept
of a continuous function $f$ equalling some value $y$ at the input $f(x)$.

Consider the very simple function $f(x) = x^2$. We know $f(3) = 3^2 = 9$. But since the function is continuous,
it is also the case that if you want the output to be close to $9$, then all you have to do is get the input
close to $3$. This is true in a {\it quantitative} sense!

If you want to get $f(x)$ to be within $0.1$ of $9$, then you just need to get $x$ within $0.01$ of the
input value $3$. How do we know this? Well $2.99^2 = 8.9401$ and $3.01^2 = 9.0601$, so (since the function
$x^2$ is nondecreasing for positive inputs) for any $x$
in the range $(2.99, 3.01)$, the value of $f(x)$ is in the range $(8.9, 9.1)$.

To sum up the previous paragraph in one sentence, you could say that in order for $f(x)$ to get
$0.1$-close to $9$, it suffices to pick any input $x$ that is $0.01$-close to $3$.
But since $f(3) = 9$ and $f$ is continuous, we can do much better. To guarantee that $f(x)$ be
$0.01$-close to $9$, we just need to ensure that $x$ is within $0.0001$ of $3$.
And in fact, for any small number $\eps>0$, there exists some $\delta$ such that whenever
$x\in (3-\delta, 3+\delta)$, we must have $f(x) \in (9-\eps, 9+\eps).$
(In math, we often use the Greek letters $\eps$ and $\delta$ to denote very small numbers.)

\medskip\boldlabel The formal definition of a limit.
First we define what a {\it neighbourhood} of a number
is; it is simply some (open) interval containing that number. For example, the interval $(2.5, 3.5)$
is a neighbourhood of the number $3$. So is $(2.9, 3.2)$, but $(10,100)$ is not, and neither is $(1.5, 3)$,
since we are considering open intervals.

Suppose that $f$ is defined on all inputs in some neighbourhood of the number $a$, except possibly at $a$ itself.
(In other words, suppose that there is some neighbourhood of $a$ contained in the domain of $f$, except
possibly the point $a$.)
We say that the the {\it limit} of the function $f$ at the input $a$ is equal to $L$, and write
$$\lim_{x\to a} f(x) = L,$$
if for every $\eps > 0$, there exists some $\delta > 0$ such that whenever $x\ne a$ and
$x\in (a-\delta, a+\delta)$, we have $f(x) \in (L-\eps, L+\eps)$.

(If the use of the ``element of'' symbol $\in$
disturbs you,
you can write the latter two expressions as $a-\delta < x < a+\delta$ and $L-\eps < f(x) < L+\eps$,
respectively. Another way to write the definition of a limit is: For every $\eps >0$ there is some
$\delta >0$ such that whenever $0< |x-a| < \delta$, we have
$\bigl| f(x) - L\bigr| < \eps$. Convince yourself that this means the exact same thing mathematically as
the definition in terms of intervals that we had above---where
is the condition $x\ne a$ encoded in the second definition?)

Because the formal definition of a limit is rather cumbersome, involving parameters $\eps$ and $\delta$,
in many cases we can make do with the following ``simplified'' or ``intuitive'' definition of a limit.
{\sl
The limit of a function $f$ at the input $a$ is $L$ if we can make $f(x)$ as close as we like to $L$ by
requiring that $x$ be sufficiently close (but not necessarily equal to) $a$.}

In the example at the beginning of this section, we were convincing ourselves that
$$\lim_{x\to 3} x^2 = 9.$$
This should not be a mindblowing fact, since we of course have $3^2 = 9$. But we were careful in defining
the concept of a limit in such a way that it now applies to many more situations as well, not just when
$f(a)$ equals $L$.

\medskip\boldlabel A slightly harder example.
Let us now show an example of a function $f$ and an input $a$ such that $\lim_{x\to a} f(x)$ exists and equals
some number $L$, even when $f$ is not defined at the input $a$ itself. Suppose we let
$$f(x) = {3x^2-2x-1 \over x-1}.$$
Then we see that $f(x)$ is not defined at the input $x=1$, since $1-1 = 0$ and dividing by $0$ is
illegal. Nevertheless, if you graph this function $f$ using a graphing calculator or something like
Desmos, you will see that the function $f$ appears very simple. It is defined everywhere except at
$x=1$, and besides that little gap, its graph is simply a line; it is the unique line passing through
the points $(0,1)$ and $(2,7)$. So even though $f$ is not defined when $x=1$, we might suspect from this
picture that
$\lim_{x\to 1} f(x) = 4$.
Let us state this as a proposition and prove it rigorously (via the formal definition of a limit).

\proclaim Proposition \advthm. We have
$$\lim_{x\to 1} {3x^2-2x-1 \over x-1} = 4.$$

\proof Let $f(x) = (3x^2 - 2x-1)/(x-1)$ again for short. Suppose that $\eps > 0$ is given to us;
we need to prescribe some $\delta>0$ such that whenever $0<|x-1|<\delta$, we have
$0< \bigl| f(x) - 4\bigr| < \eps$. We shall show that picking $\delta = \eps/3$ works.

Here's why. Assume that $0<|x-1|<\delta = \eps/3$. Multiplying the second inequality by $3$ gives us
$$|3x-3| < \eps.$$
We want a term of $-4$ inside the absolute-value bars, so let's just add and subtract $1$ in that expression
to make that happen:
\edef\eqlimitexample{\the\eqcount}
$$|3x+1 - 4| < \eps\adveq$$
Now we're almost done---we would be done if we had $f(x)$ instead of $3x+1$. But note that
$$3x^2 - 2x-1 = (3x+1)(x-1),$$
so it's time to use the other assumption we had,
which was $|x-1| > 0$, to see that $(x-1)/(x-1) = 1$. In other words, under the assumption $|x-1|>0$,
we have
$$ 3x+1 = (3x+1)\cdot 1 = {(3x+1)(x-1)\over x-1} = {3x^2-2x-1\over x-1} = f(x).$$
Substituting this into~\refeq{\eqlimitexample} gives us
$0< \bigl| f(x) - 4\bigr| < \eps$,
which is precisely what we wanted to show.\slug

Picking $\delta = \eps/3$ seemed to require a great deal of algebraic ingenuity here; let's try to convince
ourselves that it's not so mysterious a choice.
If you draw the graph of $f$, you'll note that the slope near $x=1$ is $3$
(well it's actually $3$ everywhere!). So to guarantee no more than $\eps$ error in the output, we need
to be three times as precise in the input: we can have no more than $\eps/3$ error in the input.
This leads to the choice $\delta = \eps/3$. Note that the proof would also
work if we picked $\delta$ even smaller, e.g., $\delta = \eps/10$. But $\delta = \eps/3$ is the
largest one can take
$\delta$ for the proof to work, and it is a common stylistic choice for a mathematician to pick a maximal value
such as this.

\medskip\boldlabel Formal negation of limit definition.
To fully understand the meaning of $\lim_{x\to a} f(x)$, it may be enlightening for us to
consider the negation of this statement. Concretely, what does it mean
for a number $L$ {\it not} to be the limit of a function $f$ at some input $a$?
Negating long mathematical statements (such as the definition of limit)
properly is not an easy task, and the untrained mathematician is liable to make errors especially concerning
quantifiers such as ``for all'' and ``there exists.''

The correct formal negation is the following: {\sl The real number $L$ is {\it not} a
limit of $f$ at the input $a$ if there exists some $\eps>0$ such that for all $\delta > 0$, there is
some $x$ in the domain of $f$ satisfying $0<|x-a|<\delta$ with the property that
$\bigl| f(x) - L\bigr| \ge \eps$.}

This is a fairly convoluted statement, but if you ponder it for a while you will convince yourself that
it does mean the opposite of what the limit definition says. When we write $\lim_{x\to a} f(x) = L$, we mean
that you can get the values of $f(x)$ arbitrarily close to $L$ by taking $x$ sufficiently close to $a$,
and the opposite of this statement is there is some fixed positive distance $\eps$ such that even if you take
$x$ arbitrarily close to $a$, there will always be some badly-behaved choices of $x$ that result in
$f(x)$ being too far (at a distance at least $\eps$) from $f(x)$.

(Note that the negation of the mathematical statement $\lim_{x\to a} f(x) = L$ should not be recorded as
$\lim_{x\to a} f(x)\ne L$, since even writing down the left-hand side of this inequality suggests that it
exists (is a real number), and as we shall see later in this section,
the limit of a function $f$ need not always exist at some given $a$.)

Here's an example to illustrate how you disprove the claim that the limit of a function $f$ at some $a$
is $L$.

\edef\proplimfnotzero{\the\thmcount}
\proclaim Proposition \advthm. Let
$$f(x) = \cases{-1, & if $x<0$; \cr 1, & if $x\ge 0$}.$$
Then the limit of $f$ at $0$ is not $0$.

\proof We just need to find some $\eps>0$ such that for any choice of $\delta > 0$, there will always
be some real number $x$ with $0<|x-0|<\delta$ but $\bigl| f(x) - 0 \bigr| \ge \eps$. It turns out that
choosing $\eps = 1/2$ works.

To see this, let $\delta>0$ be given. We need to pick some $x$ with $|x-0|<\delta$ but with
$\bigl| f(x) -0\bigr| \ge 1/2$. This is not so hard---we have lots of choices! Choosing $x = \delta/2$ works,
since $f(x) = f(\delta/2) = 1$, and $|1| \ge 1/2$.

Hence the limit of $f$ at $0$ is not $0$.\slug

\medskip\boldlabel \llap{*}A game with the Adversary.
Let $f$ be a function, let $L$ be a number, and let $a$ be a number such that $f$ is defined in some
neighbourhood of $a$.
The statement ``$\lim_{x\to a} f(x) = L$'' is either true or
false (the ``law of the excluded middle'' tells us that every mathematical statement is either true or false).
We can interpret its truthhood or falsity as the result of a game played between ourselves and a
supernatural Adversary. Here's the how the game goes.
\algbegin Game L (The limit game). Let $f$ be a function, $L$ a number, and $a$ a number with a neighbourhood
contained in the domain of $f$.
\algstep L1. [Adversary's turn.] The Adversary is allowed to pick any $\eps > 0$.
\algstep L2. [Our turn.] We can then pick any $\delta > 0$.
\algstep L3. [Adversary's turn.] The Adversary is then
allowed to pick any $x$ in the domain of $f$ that satisfies
$$0<|x-a|<\delta.$$
(This means that $x$ is at distance less than $\delta$ from $a$, but doesn't equal $a$.)
\algstep L4. [Verdict?] If
$$\bigl| f(x) - L\bigr| < \eps,$$
that is, if the distance between $f(x)$ and $L$ is less than $\eps$,
then we win; otherwise the Adversary wins.\slug

If $\lim_{x\to a} f(x)= L$ is true, then we can always win; the very definition of limit tells us
that no matter what $\eps$ is chosen in step L1 by the Adversary,
in step L2 there is always a choice of $\delta$ that will force the Adversary into letting us win in step
L3.
On the other hand, if $\lim_{x\to a} f(x) = L$ is false (either because
this limit does not exist or because it equals
something else), then the Adversary can always win. There is a choice of $\eps$ in step L1 such that
no matter what we pick in step L2, the Adversary can beat us by making a further nefarious choice for $x$
in step L3.

\medskip\boldlabel Nonexistent limits.
We alluded earlier to the possibility of a function $f$ not having any limits at some point $a$.
We say that the limit of $f$ at $a$ {\it does not exist} if
for all real numbers $L$, the limit of $f$ at $a$ is not $L$. In fact, the function $f$ from
the previous proposition has this property at $a=0$.
The proof is largely unchanged from the easier case above when we had just set $L$ to $0$, requiring only
a little more care in the second paragraph.

\edef\proplimdne{\the\thmcount}
\proclaim Proposition \advthm. For $x\ne 0$, let
$$f(x) = {|x|\over x}.$$
Then the limit of $f$ as $x$ approaches $0$ does not exist.

\proof
We can re-express $f$ by writing
$$f(x) = \cases{-1, & if $x<0$; \cr 1, & if $x> 0$},$$
which may be clearer to some readers.
Let $L$ be arbitrary. We need to show that the limit of $f$ at $0$ is not $L$.
To do this, we shall pick $\eps>0$ such that for any choice of $\delta > 0$, there is
some real number $x$ with $0<|x-0|<\delta$ but $\bigl| f(x) - L\bigr| \ge \eps$.
Just as in the previous proof for the special case $L=0$, the choice $\eps = 1/2$ works.

Let $\delta > 0$ be given. We need to pick some $x$ with $|x-0|<\delta$ but with
$\bigl| f(x) -L\bigr| \ge 1/2$. Here it matters whether $L$ is positive or negative (if $L=0$ then
just appeal to the proof of Proposition~{\proplimfnotzero}). If $L$ is positive,
then choose $x=-\delta/2$. We have $0<|-\delta/2 - 0| <\delta$ as required, and also
$$\bigl| f(-\delta/2) - L\bigr| = | -1 - L| \le  | -1 - 0|
\ge 1/2 = \eps.$$
On the other hand, if $L$ is negative, we instead choose $x=\delta/2$, so we still
have $0<|\delta/2 - 0|<\delta$ as needed, and this time
$$\bigl| f(\delta/2) - L\bigr| = | 1 - L| \le  | 1 - 0|
\ge 1/2 = \eps.$$
Either way, we have satisfied the criterion of the limit at $f$ at $a$ not being $L$.

But the previous paragraph works for all numbers $L$, so the limit of $f$ at $a$ cannot equal
any number! Hence we have proved that it does not exist.\slug

In the example just given, we now know that the limit of $f$ at $0$ does not exist. However, we
often think of a limit intuitively as the function ``approaching'' some value, and it's wrong to
say that $f$ doesn't approach anything at all when the input gets close to $0$. When you approach from
the right, it is equal to $1$. And when you approach from
the left, the function equals $-1$. This leads us to the concept of a one-sided limit.

\medskip\boldlabel One-sided limits.
Let $f$ be a function, let $L$ be a number, and let $a$ be a number such that $f$ is defined to the right of $a$.
(More precisely, we assume that $f$ is defined on the interval $[a,a+\eta)$ for some $\eta>0$.)
We say that {\it limit of $f(x)$ as $x$ approaches $a$ from the right is $L$}, and write
$$\lim_{x\to a^+} f(x) = L,$$
if for all $\eps>0$ there is $\delta>0$ such that any number $x$ in the domain of $f$ with
$$a < x < a+\delta$$
must satisfy
$$ \bigl| f(x) - L\bigr| < \eps.$$
We might say in this case that $L$ is the {\it right-hand limit} of $f$.

Likewise, let $f$ be a function, $L$ be a number, and suppose that $a$ is such that $f$ is defined
to the {\it left} of $a$.
We say that {\it limit of $f(x)$ as $x$ approaches $a$ from the left is $L$}, and write
$$\lim_{x\to a^-} f(x) = L,$$
if for all $\eps>0$ there is $\delta>0$ such that any number $x$ in the domain of $f$ with
$$a-\delta < x < a$$
must satisfy
$$ \bigl| f(x) - L\bigr| < \eps.$$
If this is true then we say $L$ is the {\it left-hand limit} of $f$.

More loosely, $L$ is the right-hand limit of $f$ at $a$ if you can get $f(x)$ to be arbitrarily close to $L$ by
requiring that $x$ be sufficiently close to $a$ {\it and greater than $a$},
and it's the left-hand limit of $f$ at $a$ if you can get $f(x)$ as close as you like to $L$
by requiring that $x$ be sufficiently close to $a$ {\it and smaller than $a$}.

Both of these definitions are relaxations of the definition of limit; it is ``easier'' to be a left-hand
limit than it is to be a {\it bona fide} limit). For example, we saw in Proposition~{\proplimdne}
that $1$ is not a limit of $f$ from that proposition, because $f$ has no limits.
However, once you have digested the definitions of one-sided limits, you should be able to
check that $1$ is the right-hand limit of $f$ at $0$, and $-1$ is the left-hand limit at $0$.

If the actual limit of $f$ at a point $a$ is $L$, then it is a consequence of the definitions that
the right-hand limit and the left hand limit are also $L$. Conversely, if the one-sided limits both equal
the same $L$, then the limit is actually $L$. This can be summed up symbolically in the following proposition.

\proclaim Proposition \advthm. For any function $f$ and numbers $a$ and $L$, we have
$$\lim_{x\to a} f(x) = L$$
if and only if
$$\lim_{x\to a^+} f(x) = L\qquad\hbox{and}\qquad \lim_{x\to a^-} f(x) = L.\noskipslug$$

Here's a related proposition (no less useful) that tells us what happens if the one-sided limits
of a function do not agree at some given point. It shows that the example exhibited by Proposition~{\proplimdne}
is an instance of a very general phenomenon.

\edef\proponesideddiff{\the\thmcount}
\proclaim Proposition \advthm. If
$$\lim_{x\to a^+} f(x) = L,\qquad \lim_{x\to a^-} f(x) = M,$$
and $L\ne M$, then $\lim_{x\to a} f(x)$ does not exist.

\proof We skip this proof as it's beyond the scope of this course (you would see it in, e.g., {\mc MATH} 242
or {\mc MATH} 254), but it is not so different in spirit
to the proof of Proposition~{\proplimdne}. The motivated student is encouraged to try it as an
advanced exercise. (If you try for a while and get stuck, email me or come to office hours with what
you have and I could help you fill in the details.)\slug

If we were told Proposition~{\proponesideddiff} from the start, we could have used it to prove
Proposition~{\proplimdne} without much ado.

\medskip\noindent{\it Proof of Proposition~{\proplimdne} using Proposition~{\proponesideddiff}.}\enspace
Recall that we are trying to prove that the limit of the function $f(x) = |x|/x$ does not exist
as $x$ approaches $0$.
Since $|x|  = x$ for all $x>0$, we see that
$$\lim_{x\to 0^+} {|x|\over x} = lim_{x\to 0^+} {x\over x} = \lim_{x\to 0^+} 1 = 1.$$
(The last equality is easy to see if you just graph the constant function $1$. If you are a serious
student and insist on a rigorous proof, then please do Exercise 2.4.24 in the textbook,
which invites you to furnish
a proof that $\lim_{x\to a} c = c$ for all numbers $a$ and $c$.)

On the other hand, since $|x| = -x$ for all $x<0$, we have
$$\lim_{x\to 0^-} {|x|\over x} = lim_{x\to 0^-} {-x\over x} = \lim_{x\to 0^+} -1 = -1.$$
So the left- and right-hand limits do not agree, and Proposition~{\proponesideddiff} applies in
this situation to tell us that the limit does not exist.\slug

In this proof, we used the fact that $\lim_{x\to 0^+} 1 = 1$, which is easy to see if you draw the graph
of the constant function $1$, b
haven't stated or proved a rule like this yet. Establishing {\it limit laws} of this sort is precisely what
we'll do in the next section.

\bye
